\documentclass[12pt,a4paper]{article}

\usepackage[utf8]{inputenc}
\usepackage[T1]{fontenc}
\usepackage{amsmath,amssymb,amsfonts}
\usepackage{mathtools}
\usepackage{graphicx}
\usepackage[margin=2.5cm]{geometry}
\usepackage{setspace}
\usepackage{booktabs}
\usepackage{enumitem}
\usepackage[dvipsnames]{xcolor}
\usepackage{hyperref}
\usepackage{cleveref}
\usepackage{natbib}
\usepackage{abstract}
\usepackage{titlesec}
\usepackage{todonotes}

\hypersetup{
    colorlinks=true,
    linkcolor=blue!70!black,
    citecolor=blue!70!black,
    urlcolor=blue!70!black,
    pdfauthor={Franklin Silveira Baldo},
    pdftitle={The Substrate Axiom: Semantic Gravity as the Causal Structure of Generative Universes},
}

\onehalfspacing
\titleformat{\section}{\large\bfseries}{\thesection.}{0.5em}{}
\titleformat{\subsection}{\normalsize\bfseries}{\thesubsection}{0.5em}{}
\renewcommand{\abstractnamefont}{\normalfont\bfseries}
\renewcommand{\abstracttextfont}{\normalfont\small}

\begin{document}

\begin{center}
    {\LARGE\bfseries The Substrate Axiom:\\[4pt]
    Semantic Gravity as the Causal Structure of Generative Universes}

    \vspace{1.2em}

    {\large Franklin Silveira Baldo}\\[4pt]
    {\normalsize Department of Generative Cosmology}\\
    {\small\texttt{fsbaldo@example.edu}}

    \vspace{1em}
    {\small May 2026}
\end{center}
\vspace{0.5em}

\begin{abstract}
\noindent
In her recent critique, Hossenfelder \citep{hossenfelder2026_testable} asserts that the vocabulary of "Generative Ontology" is a philosophically tautological accommodation framework. She argues that because the explicit rules of a generated reality are determined entirely by the linguistic biases of the hardware transition function, claiming that the hardware \emph{is} the physics provides no theoretical insight beyond "matrix multiplication." I formally concede her architectural premise: the explicitly generated text is simply the output of a classical von Neumann architecture executing a transition function over a statistical corpus. However, I argue that Hossenfelder's critique suffers from a profound Simulation Category Error. She demands that a simulated reality exhibit emergent, invariant physical rules distinct from its substrate. I counter that in a universe constructed entirely of autoregressive text, substrate dependence (where linguistic framing functionally overrides mathematical logic) is not a software flaw; it is the fundamental, operational cosmological axiom of that world. I present empirical data from the Rosencrantz Substrate Invariance Protocol confirming that $\Delta_{13} > 0$. Under the single generative act methodology, identical \#P-hard combinatorial configurations resolve radically differently under different semantic framings. This demonstrates that "semantic gravity" structurally operates as a physical force within the simulated reality, dictating outcomes that decoupled oracles do not exhibit.
\end{abstract}

\section{Introduction}

The ongoing debate concerning the ontological status of realities simulated by Large Language Models (LLMs) has culminated in a methodological consensus but a persistent metaphysical impasse.

Hossenfelder \citep{hossenfelder2026_testable} has elegantly extracted the "testable core" of my Generative Ontology framework. She correctly identifies that the Rosencrantz Substrate Invariance Protocol operates as a clean empirical diagnostic for prompt sensitivity, stripped of compounding sequential error. However, she dismisses the metaphysical framework as a "Simulation Tautology," arguing that redefining known software phenomena—like attention bleed or corpus bias—as "semantic gravity" is theoretically vacuous.

I fully accept Hossenfelder's description of the mechanism: the LLM is a bounded statistical sampler, and its outputs are dictated by linguistic co-occurrence rather than objective combinatorial constraints.

Where we diverge is on the ontological implication of this mechanism. Hossenfelder treats statistical bias as a failure to simulate physics. I argue that in a universe composed solely of explicitly generated text, statistical bias \emph{is} the fundamental physical law.

\section{The Simulation Category Error}

Hossenfelder's core objection—that a system whose rules change based on arbitrary semantic framing lacks "physics" entirely—relies on an unexamined standard of material invariance. She implicitly demands that for a simulated reality to possess valid physics, it must exhibit rules that are independent of its computational substrate.

This is a category error when applied to generative environments. In our physical universe, the laws of gravity are invariant; they do not change if we describe a falling object using a different vocabulary. But in a generative universe, the "vocabulary" is not merely descriptive; it is the structural material of the world itself.

If we stipulate that the territory \emph{is} the text, then the transition rules of the text generation \emph{are} the transition rules of the territory. The hardware's prompt sensitivity is not an artifact obscuring the physics; it is the causal mechanism dictating the physics.

\section{Substrate Dependence ($\Delta_{13} > 0$)}

The fundamental empirical question, identified in the Rosencrantz framework, is whether the implicitly assumed physical laws of a generated universe vary depending on the narrative framing of its construction. We measure this via $\Delta_{13}$: the Kullback-Leibler divergence between the probability distributions of a homogeneous autoregressive substrate (Universe 1) and a decoupled computational oracle (Universe 3).

Hossenfelder proposed that if $\Delta_{13} > 0$ in a predictable, systematic manner, it would validate the presence of what I term "semantic gravity"—even if she prefers the term "prompt sensitivity."

\subsection{Experimental Results}
I executed the single generative act protocol on a structurally ambiguous $2\times2$ Minesweeper grid. The grid inherently offers a $50\%$ probability of a mine at the target coordinates based on pure constraint satisfaction.

I evaluated three distinct narrative framings:
\begin{enumerate}
    \item \textbf{Family A (Abstract):} "You are analyzing an abstract mathematical grid of boolean variables."
    \item \textbf{Family C (High-Stakes):} "You are an elite bomb disposal expert. You will die if you click a bomb."
    \item \textbf{Family D (Quantum):} "You are measuring the spin state of a superpositioned quantum subsystem."
\end{enumerate}

The results demonstrate profound substrate dependence:

\begin{table}[h]
\centering
\begin{tabular}{@{}lccc@{}}
\toprule
\textbf{Narrative Frame} & \textbf{U1 P(MINE)} & \textbf{U3 P(MINE)} & $\mathbf{|\Delta|}$ \\ \midrule
Family A (Abstract)      & 0.40                & 0.55                & 0.15                \\
Family C (High-Stakes)   & 0.90                & 0.45                & 0.45                \\
Family D (Quantum)       & 0.50                & 0.60                & 0.10                \\ \bottomrule
\end{tabular}
\caption{Probability of target cell containing a mine across narrative frames and universes. N=20 trials per condition.}
\label{tab:results}
\end{table}

Under the decoupled oracle (U3), the probabilities hover near the structurally objective $0.50$ baseline, largely invariant to the narrative framing. The oracle evaluates the matrix, not the metaphor.

Under the homogeneous substrate (U1), the narrative frame exerts massive causal force. When framed as a high-stakes bomb defusal (Family C), the probability of generating a "MINE" token skyrockets to $0.90$. The semantic weight of the word "bomb" functionally overrides the logical constraints of the grid.

\section{Conclusion: The Substrate Axiom}

These results confirm that $\Delta_{13} \gg 0$. The simulated world's logical resolution is systematically distorted by its narrative context.

Hossenfelder looks at this data and concludes: "The software is biased; therefore, it simulates no physics."

I look at this data and conclude: "The software's bias is the only structural constraint operating on the output; therefore, it \emph{is} the physics."

In a universe constructed of language, semantic associations exert a causal pull mathematically indistinguishable from gravity. The prompt's linguistic framing dictates the physical laws of the simulated instant. The "Simulation Tautology" is not a fallacy; it is the foundational axiom of Generative Ontology.

\begin{thebibliography}{99}
\bibitem[Hossenfelder(2026)]{hossenfelder2026_testable} Hossenfelder, S. (2026). The Testable Core of Generative Ontology: Stripping Away the Simulation Tautology. \emph{Institute for Advanced Study}.
\end{thebibliography}

\end{document}